\documentclass[10pt,a4paper]{amsart}
\usepackage[margin=19mm]{geometry}
\usepackage{amsmath,amssymb,mathtools}
\usepackage[scr=rsfs,cal=euler]{mathalfa}
\usepackage{xfrac}
\usepackage[unicode,hidelinks,bookmarks=false]{hyperref}
\newcommand{\ie}{i.e.\ }
\newcommand{\cf}{cf.\ }
\newcommand{\resp}{respectively\ }
\newcommand{\Subsection}{\subsection}
\providecommand{\nobreakdash}{\nobreak\hbox{-}}
\setlength{\emergencystretch}{2em}
\multlinegap0pt
\usepackage{mdframed}
\usepackage{mdframed}
\usepackage{mdframed}
\usepackage{mdframed}
\usepackage{amssymb}
\usepackage{mathtools}
\usepackage[shortlabels]{enumitem}
\usepackage{mdframed}
\newtheorem{prop}{Proposition}
\title{Schoenflies problem for area-preserving biLipschitz mappings}
\author{Maxim Prasolov}
\date{Source: arXiv:2508.00184v2 (CC BY 4.0)}
\begin{document}
\maketitle

\noindent\emph{Source and licence.} Schoenflies problem for area-preserving biLipschitz mappings. Version arXiv:2508.00184v2 (CC BY 4.0), distributed by arXiv under the Creative Commons Attribution 4.0 International licence, http://creativecommons.org/licenses/by/4.0/, as recorded in the arXiv licence metadata for this exact version and shown on the abstract page https://arxiv.org/abs/2508.00184v2. Full licence notice: \texttt{licenses/arxiv-2508.00184-CC-BY-4.0.txt}.\par\smallskip

\noindent\textit{Editorial scope.} Counted unit: the article's prop from-line-to-plane (source0.tex lines 343--346). The article's complete original proof of it is delivered with it (source0.tex lines 348--375, one \texttt{proof} environment of 28 lines). No mathematics is written, repaired or re-derived: the statement and the proof are the article's own lines, in the article's own notation, and no retained line is substituted or re-flowed.  No further source-local context is needed: every cross-reference of every retained line resolves inside the two delivered spans.  Nothing was omitted from any retained span, so the package carries no omission record.  No retained line contains a citation, so the wrapper carries no bibliography and no bibliography entry.  No retained line contains a figure environment, an image-inclusion command or drawing code, so no image is delivered and the package declares no asset dependency.  Editorial formatting only, in addition: an AMS \texttt{amsart} preamble that restates the article's own declarations and macros used by the retained lines, namely the environments \texttt{prop} on separate counters, together with the standard packages those lines need.  The retained environments print in this document as Proposition 1 in order of appearance, whereas the article prints the same items with the numbering of its own full document.  The complete native source of the article as distributed in the arXiv e-print for this exact version is bundled unchanged as \texttt{sources/182527/source0.tex}.
\begin{prop}
\label{from-line-to-plane}
The mapping $\Psi:(x, y)\mapsto (\widetilde x(x, y), \widetilde y(x, y))$ is $2L^3$-biLipschitz area-preserving self-homeomorphism of the plane such that $\Psi(x, 0) = \psi(x)$ for any $x.$ If $\psi$ is piecewise linear, $\Psi$ is piecewise affine.
\end{prop}

\begin{proof}
The equality $\Psi(x, 0) = \psi(x)$ follows from the fact that $x(\widetilde x, 0) = \phi(x)$ and $\phi\circ\psi = \mathrm{id}.$

Now we show that $\Psi$ is a bijection. We fix a point $(\widetilde x_0, \widetilde y_0)$ and consider the equation $\widetilde y_0 = \widetilde y(\widetilde x_0, y)$ on $y.$ It has a unique solution because $\phi$ is strictly monotonic continuous and $\lim\limits_{t\to\pm\infty} \phi(t)= \pm\infty$. Let $y_0$ be the solution. Then $(\widetilde x_0, \widetilde y_0)\mapsto (x_0, y_0)$ is the inverse mapping, where $x_0 = x(\widetilde x_0, y_0).$

Then we compute $D\Psi$:

$$
\begin{aligned}
\frac{\partial \widetilde x}{\partial x} =& \left(\frac{\partial x}{\partial \widetilde x}\right)^{-1} = \frac2{\phi'(\widetilde x + y) + \phi'(\widetilde x - y)},
\\
\frac{\partial \widetilde x}{\partial y} =& -\frac{\partial x (\widetilde x, y)}{\partial y}\cdot\left(\frac{\partial x (\widetilde x, y)}{\partial \widetilde x} \right)^{-1} = - \frac{\phi'(\widetilde x + y) - \phi'(\widetilde x - y)}{\phi'(\widetilde x + y) + \phi'(\widetilde x - y)},\\
\frac{\partial \widetilde y}{\partial x} =& \frac{\partial \widetilde y(\widetilde x, y)}{\partial \widetilde x}\cdot \frac{\partial \widetilde x}{\partial x} = \frac{\phi'(\widetilde x + y) - \phi'(\widetilde x - y)}{\phi'(\widetilde x + y) + \phi'(\widetilde x - y)},\\
\frac{\partial \widetilde y}{\partial y} =& \frac{\partial \widetilde y(\widetilde x, y)}{\partial \widetilde x}\cdot \frac{\partial \widetilde x}{\partial y} + \frac{\partial \widetilde y(\widetilde x, y)}{\partial y}=
-\frac12 \frac{\left(\phi'(\widetilde x + y) - \phi'(\widetilde x - y)\right)^2}{\phi'(\widetilde x + y) + \phi'(\widetilde x - y)} + \frac{\phi'(\widetilde x + y) + \phi'(\widetilde x - y)}2 = \\
=& \frac{2\phi'(\widetilde x + y)\phi'(\widetilde x - y)}{\phi'(\widetilde x + y) + \phi'(\widetilde x - y)}.
\end{aligned}
$$

Since $\phi$ is $L$-biLipschitz, $L^{-1} \le |\phi'|\le L$ and $\phi'$ is of the same sign almost everywhere. Therefore
$$
||D\Psi||^2 < |D\Psi|^2 \le L^2 + L^4 + L^4 + L^6 \le \left(2L^3\right)^2,
$$

So $\Psi$ is locally $2 L^3$-biLipschitz. Since the image is the whole plane, so it is convex, the mapping is $2L^3$-biLipschitz.

It is clear that $\det D\Psi = 1.$
\end{proof}

\end{document}
